% Switch-on surge (rectifier-surge): the set-up a person would wire, semantically the deck's.
%   the regulation deck with every capacitor discharged at t = 0 (uic), the cathode already hot; read the plate currents over the first two mains periods
\documentclass[tikz, border=10pt]{standalone}
\usepackage{bench}
\begin{document}
\begin{tikzpicture}
\benchtitle{0,8.4}{Switch-on surge}
\def\gnd{0.8}
\coordinate (pt) at (1.0,5.0);
\coordinate (pb) at (1.0,1.6);
\draw[set] (pt) to[sV, l_={$\sqrt 2\,V_{ac}\sin(\omega t + \phi)$}] (1.0,3.3) coordinate (ctap);
\draw[set] (ctap) to[sV, l_={antiphase}] (pb);
\draw (ctap) -- ++(-0.6,0) node[ground, rotate=-90]{};
\draw (pt) to[R, l={$R_s$}] ++(2.2,0) coordinate (pa);
\coordinate (pam) at ($(pa)+(1.6,0)$);
\ammeterh{pa}{pam}{$i_{a1}(t)$}
\draw (pb) to[R, l_={$R_s$}] ++(2.2,0) coordinate (pbx);
\coordinate (pbm) at ($(pbx)+(1.6,0)$);
\ammeterh{pbx}{pbm}{$i_{a2}(t)$}
\node[diodetube, rotate=90] (D1) at ($(pam)+(1.3,0)$) {};
\node[diodetube, rotate=90] (D2) at ($(pbm)+(1.3,0)$) {};
\draw (pam) -- (D1.anode); \draw (pbm) -- (D2.anode);
\draw (D1.cathode) -- ++(0.5,0) coordinate (k1); \draw (D2.cathode) -- ++(0.5,0) coordinate (k2);
\draw (k1) -- (k2);
\coordinate (c) at ($(k1)!0.5!(k2)+(1.0,0)$);
\draw ($(k1)!0.5!(k2)$) -- (c);
\draw (c) to[C, l={$C_1$, discharged}] (c |- 0,\gnd) node[ground]{};
\draw (c) -- ++(3.4,0) coordinate (il);
\draw[set] (il) to[I, l={$I_{dc}$}] (il |- 0,\gnd) node[ground, black]{};
\node[font=\small, align=center] at ($(D1)!0.5!(D2)$) {DUT\\(twin diode)};
% the definitions
\definitions{0,-0.1}{%
  $\phi = 0$ and $90^\circ$ (switched at the zero crossing and at the peak);\quad $C_1$ and $R_s$ as in the regulation deck;\quad $V_{ac}$ per plate as the sheets' rating charts\\[2pt]
  read the first peak plate current per plate, beside the hot-switching ratings (2.5 A for the 5Y3, 1 A for the 6X4) and the minimum $R_s$ of GE's Rating Chart III\\[2pt]
  what it cannot show: a directly heated cathode's warm-up, which delays and softens the surge (no model has a heater)%ENDDEFS
}
\end{tikzpicture}
\end{document}
